\begin{explanationbox}
	\textbf{\textcolor{CExplanation}{一句话直觉}}
	一次分式函数通过分离常数可化为“常数 + 反比例”形式，值域由常数项唯一决定，图像由反比例函数平移得到。

	\medskip
	\textbf{\textcolor{CExplanation}{核心拆解}}
	\begin{description}[style=nextline, leftmargin=2.8em, labelsep=0.5em, itemsep=0.45em]
		\item[\textbf{要点一（分离常数）：}]
			利用多项式除法或配拆法，将分子中含 $x$ 项系数与分母保持一致，从而得到一项常数和一项含分母的一次式。
		\item[\textbf{要点二（值域推导）：}]
			常数项 $a/c$ 是反比例部分的垂直平移量；反比例部分可取任意非零实数，故原函数值域排除 $a/c$。
	\end{description}

	\medskip
	\textbf{\textcolor{CExplanation}{几何本质（双曲线）}}
	一次分式函数 $y=\dfrac{ax+b}{{cx+d}}$ 的图像是等轴双曲线，其对称中心为 $\left(-\dfrac{d}{c},\dfrac{a}{c}\right)$，渐近线为 $x=-\dfrac{d}{c}$ 和 $y=\dfrac{a}{c}$。

	\medskip
	\textbf{\textcolor{CExplanation}{代数意义（恒等变形）}}
	分离常数实现了将分式函数拆分为多项式函数与真分式之和，便于研究值域、单调性及图像变换。

	\medskip
	\textbf{\textcolor{CExplanation}{考点价值（常考题型）}}
	\begin{itemize}[leftmargin=*, itemsep=0.5em, topsep=0.4em]
		\item \textbf{考法一（值域问题）：} 直接套用值域公式 $y\neq a/c$，避免分类讨论。
		\item \textbf{考法二（对称性质）：} 由 $(-d/c,a/c)$ 直接得对称中心，用于图像变换或求单调区间。
	\end{itemize}

	\medskip
	\textbf{\textcolor{CExplanation}{顿悟点}}
	分离常数后，看到 $y-\dfrac{a}{c} = \dfrac{bc-ad}{c}\cdot\dfrac{1}{{c(x+\frac{d}{c})}}$，本质上就是反比例函数 $Y=\dfrac{k}{X}$，其中 $Y=y-\dfrac{a}{c}$，$X=x+\dfrac{d}{c}$，$k=\dfrac{bc-ad}{{c^2}}$。因此所有性质可直接沿用反比例函数结论。

	\medskip
	\textbf{\textcolor{CExplanation}{使用场景}}
	\begin{itemize}[leftmargin=*, itemsep=0.5em, topsep=0.4em]
		\item \textbf{场景一（求值域）：} 已知一次分式函数，直接得值域。
		\item \textbf{场景二（图像变换）：} 求对称中心、渐近线或平移关系。
	\end{itemize}

\end{explanationbox}
